\documentclass[12pt, a4paper]{article}

% --- Paquetes de Idioma y Codificación ---
\usepackage[utf8]{inputenc}
\usepackage[T1]{fontenc}
\usepackage[spanish]{babel}

% --- Paquetes Matemáticos ---
\usepackage{amsmath}
\usepackage{amssymb}
\usepackage{amsfonts}
\usepackage{mathtools}

% --- Paquetes de Diseño y Formato ---
\usepackage{geometry}
\geometry{margin=2.5cm}
\usepackage{xcolor}
\usepackage{titlesec}
\usepackage{fancyhdr}
\usepackage{tcolorbox}
\tcbuselibrary{theorems, skins, breakable}
\usepackage{booktabs}
\usepackage{array}
\usepackage{hyperref}
\usepackage{enumitem}
\usepackage{graphicx}

% --- Colores Personalizados GEM ---
\definecolor{gemblue}{RGB}{0, 51, 153}
\definecolor{gemgold}{RGB}{204, 153, 0}
\definecolor{gemdark}{RGB}{20, 20, 20}
\definecolor{cosmicpurple}{RGB}{102, 0, 153}
\definecolor{alertred}{RGB}{200, 0, 0}

% --- Configuración de Encabezados y Pies ---
\pagestyle{fancy}
\fancyhf{}
\fancyhead[L]{\textcolor{cosmicpurple}{\textsc{Protocolo GEM-01}}}
\fancyhead[R]{\textsc{Validación Experimental}}
\fancyfoot[C]{\textcolor{gray}{\thepage}}
\renewcommand{\headrulewidth}{0.4pt}

% --- Estilos de Título ---
\titleformat{\section}{\Large\bfseries\color{cosmicpurple}}{\thesection}{1em}{}[\titlerule]
\titleformat{\subsection}{\large\bfseries\color{gemdark}}{\thesubsection}{1em}{}

% --- Título del Documento ---
\title{\textbf{\Huge \textcolor{cosmicpurple}{Protocolo de Validación Experimental GEM-01}} \\[0.5cm]
\Large \textit{Medición de Reducción de Inercia Topológica en Agua bajo Resonancia Escalar} \\[0.3cm]
\large \textcolor{gemgold}{\textbf{Ciencia Abierta / Replicable}}}
\author{Investigación Colaborativa G.E.M. \\ \small Basado en el Teorema 02: La Válvula de Vacío \\ \small Junio 2026}
\date{}

\begin{document}

\maketitle
\thispagestyle{empty}
\newpage

\tableofcontents
\newpage

% ==========================================
% SECCIÓN 1: OBJETIVO Y FUNDAMENTO TEÓRICO
% ==========================================
\section{Objetivo y Fundamento Teórico}

\subsection{Objetivo Principal}
Validar experimentalmente el \textbf{Teorema 02 del Modelo GEM} (La Válvula de Vacío) mediante la medición de reducción de inercia topológica en una muestra de agua ultrapura sometida a un gradiente de potencial escalar $\nabla w$ generado por nulificación vectorial a la frecuencia de resonancia de $16.2$ Hz.

\subsection{Hipótesis a Validar}
Según el Teorema 02, al aplicar un gradiente topológico $\nabla w$ sobre un transductor geométrico (agua con ángulo molecular $105^\circ$), se extrae energía del vacío cuántico según:
\begin{equation}
\dot{E}_{ext} = \eta_{GEM} \cdot \oint_{\partial V} [\nabla w \cdot \mathbf{J}_g] d\mathbf{S}
\end{equation}

Esta energía extraída se manifiesta como una \textbf{reducción de inercia} $\Delta I$ proporcional a:
\begin{equation}
\Delta I = \frac{E_{ext}}{1836 \cdot C_{105}}
\end{equation}

Donde:
\begin{itemize}
    \item $1836$ es la relación de masa protón/electrón (factor de confinamiento topológico)
    \item $C_{105} \approx 0.98$ es el factor de corrección geométrica del agua
    \item $\eta_{GEM} = \frac{\phi}{\pi} \cdot \delta^2 \approx 2.06 \times 10^{-4}$ es la eficiencia de acoplamiento
\end{itemize}

\subsection{Predicción Cuantitativa}
Para una muestra de $50$ mL de agua Milli-Q sometida a un gradiente $\nabla w$ generado por bobinas de Helmholtz a $16.2$ Hz con amplitud de $5$ Vpp, el modelo predice:
\begin{itemize}
    \item \textbf{Reducción de masa esperada:} $\Delta m \approx 0.005$ g a $0.05$ g ($0.01\%$ - $0.1\%$)
    \item \textbf{Tiempo de estabilización:} $5$ - $15$ minutos
    \item \textbf{Histéresis topológica:} La reducción persiste $5$-$10$ minutos tras apagar el campo
\end{itemize}

% ==========================================
% SECCIÓN 2: MATERIALES Y EQUIPAMIENTO
% ==========================================
\section{Materiales y Equipamiento Requerido}

\subsection{Componentes del Rig GEM-01}

\begin{table}[h!]
\centering
\caption{Lista de Materiales del Rig GEM-01}
\begin{tabular}{@{}p{6cm}p{4cm}p{5cm}@{}}
\toprule
\textbf{Componente} & \textbf{Especificación} & \textbf{Función GEM} \\ \midrule
Bobinas de Helmholtz (3 ejes) & Diámetro 20 cm, 50 espiras, AWG 22 & Generar nulificación vectorial $\mathbf{v}=0$ \\ \addlinespace
Generador de señales DDS & Estabilidad 0.001 Hz, salida 0-10 Vpp & Generar señal de 16.2 Hz pura \\ \addlinespace
Amplificador de audio & 20 W, respuesta 10 Hz-100 kHz & Amplificar señal para bobinas \\ \addlinespace
Balanza analítica & Precisión 0.1 mg, capacidad 200 g & Medir reducción de masa $\Delta m$ \\ \addlinespace
Vaso de precipitados & Cuarzo o vidrio borosilicato, 100 mL & Contener muestra de agua \\ \addlinespace
Agua Milli-Q & Resistividad 18.2 MΩ·cm, TOC < 5 ppb & Transductor geométrico ($105^\circ$) \\ \addlinespace
Termistor NTC & Precisión 0.01°C, respuesta rápida & Monitorear $\Delta T$ (entropía) \\ \addlinespace
Jaula de Faraday & Malla cobre + blindaje Mu-Metal & Aislar de EMI externo \\ \addlinespace
Osciloscopio & 100 MHz, 4 canales & Verificar nulificación vectorial \\ \addlinespace
Gaussímetro & Rango 0-200 mT, resolución 0.01 mT & Medir campo magnético residual \\ \bottomrule
\end{tabular}
\end{table}

\subsection{Configuración del Sistema}

\begin{figure}[h!]
\centering
\fbox{\parbox[c][6cm][c]{12cm}{\centering \textbf{DIAGRAMA ESQUEMÁTICO DEL RIG GEM-01} \\[1cm]
[Espacio para insertar diagrama SVG/PDF] \\[1cm]
Generador DDS $\rightarrow$ Amplificador $\rightarrow$ Bobinas Helmholtz (X,Y,Z) \\
$\downarrow$ \\
Muestra de Agua (50 mL) en Balanza \\
$\downarrow$ \\
Termistor NTC + Osciloscopio}}
\caption{Configuración experimental del Rig GEM-01}
\end{figure}

% ==========================================
% SECCIÓN 3: PROCEDIMIENTO EXPERIMENTAL
% ==========================================
\section{Procedimiento Experimental Paso a Paso}

\subsection{Fase 1: Preparación y Calibración (Día 1)}

\begin{enumerate}[label=\textbf{Paso \arabic*:}, leftmargin=*]
    \item \textbf{Limpieza del Sistema:}
    \begin{itemize}
        \item Limpiar el vaso de precipitados con ácido nítrico al 10\% y enjuagar 3 veces con agua Milli-Q.
        \item Secar en estufa a 60°C durante 30 minutos.
        \item Manipular siempre con guantes de nitrilo para evitar contaminación.
    \end{itemize}
    
    \item \textbf{Calibración de la Balanza:}
    \begin{itemize}
        \item Nivelar la balanza usando el nivel de burbuja integrado.
        \item Calibrar con pesas patrón certificadas (10 g, 50 g, 100 g).
        \item Verificar deriva térmica: la lectura debe ser estable ($\pm 0.1$ mg) durante 30 minutos.
        \item Registrar temperatura ambiente: debe ser $20 \pm 0.5^\circ$C.
    \end{itemize}
    
    \item \textbf{Configuración de Bobinas de Helmholtz:}
    \begin{itemize}
        \item Montar las 3 bobinas en configuración ortogonal (X, Y, Z).
        \item Conectar en serie con polaridad correcta (verificar con multímetro).
        \item Medir resistencia total: debe ser $R \approx 10-15 \, \Omega$.
        \item Verificar inductancia: $L \approx 5-10$ mH.
    \end{itemize}
    
    \item \textbf{Calibración del Generador DDS:}
    \begin{itemize}
        \item Configurar frecuencia: $f = 16.200$ Hz (exacto).
        \item Configurar amplitud: $V_{pp} = 5.00$ V.
        \item Configurar forma de onda: Senoidal pura (THD < 0.1\%).
        \item Verificar con osciloscopio: medir frecuencia real y amplitud.
    \end{itemize}
\end{enumerate}

\subsection{Fase 2: Línea Base y Control (Día 2)}

\begin{enumerate}[label=\textbf{Paso \arabic*:}, leftmargin=*, start=5]
    \item \textbf{Medición de Línea Base (Sin Campo):}
    \begin{itemize}
        \item Llenar el vaso con exactamente $50.00$ mL de agua Milli-Q (usar pipeta volumétrica).
        \item Colocar el vaso en el centro de las bobinas de Helmholtz (intersección de los 3 ejes).
        \item Colocar el termistor NTC en contacto con la pared exterior del vaso.
        \item Cerrar la jaula de Faraday.
        \item Esperar 30 minutos para estabilización térmica.
        \item Registrar masa inicial $m_0$ cada minuto durante 10 minutos.
        \item Calcular promedio: $\bar{m}_0 = \frac{1}{10} \sum_{i=1}^{10} m_i$
        \item Calcular desviación estándar: $\sigma_0$ (debe ser $< 0.2$ mg).
    \end{itemize}
    
    \item \textbf{Verificación de Nulificación Vectorial:}
    \begin{itemize}
        \item Encender el generador DDS a 16.2 Hz, 5 Vpp.
        \item Conectar el gaussímetro en el centro de las bobinas.
        \item Ajustar fase y amplitud de cada bobina (X, Y, Z) hasta minimizar el campo magnético.
        \item \textbf{Criterio de éxito:} $|\mathbf{B}| < 0.05$ mT en los 3 ejes.
        \item Verificar con osciloscopio que los campos $\mathbf{E}$ y $\mathbf{B}$ están en oposición de fase ($180^\circ$).
    \end{itemize}
\end{enumerate}

\subsection{Fase 3: Experimento Principal (Día 3)}

\begin{enumerate}[label=\textbf{Paso \arabic*:}, leftmargin=*, start=7]
    \item \textbf{Aplicación del Gradiente $\nabla w$:}
    \begin{itemize}
        \item Verificar que la muestra está en las mismas condiciones que la línea base.
        \item Registrar masa inicial $m_{inicio}$.
        \item Registrar temperatura inicial $T_{inicio}$.
        \item \textbf{Encender el generador DDS} a $t=0$.
        \item Iniciar registro automático de datos cada 10 segundos:
        \begin{itemize}
            \item Masa $m(t)$
            \item Temperatura $T(t)$
            \item Voltaje en bobinas $V(t)$
            \item Campo magnético residual $B(t)$
        \end{itemize}
    \end{itemize}
    
    \item \textbf{Monitoreo en Tiempo Real (60 minutos):}
    \begin{itemize}
        \item Minuto 0-10: Fase de transición (posibles oscilaciones).
        \item Minuto 10-50: Fase de estabilización (buscar meseta).
        \item Minuto 50-60: Fase de medición estable.
        \item \textbf{Alerta:} Si $\Delta T > +2.0^\circ$C, APAGAR INMEDIATAMENTE (fuga vectorial).
    \end{itemize}
    
    \item \textbf{Apagado y Medición de Histéresis:}
    \begin{itemize}
        \item En $t=60$ min, \textbf{apagar el generador DDS}.
        \item Continuar registrando datos cada 10 segundos durante 30 minutos adicionales.
        \item Esto mide la \textbf{histéresis topológica}: cuánto persiste el efecto tras apagar el campo.
    \end{itemize}
\end{enumerate}

\subsection{Fase 4: Análisis de Datos}

\begin{enumerate}[label=\textbf{Paso \arabic*:}, leftmargin=*, start=10]
    \item \textbf{Cálculo de la Reducción de Masa:}
    \begin{equation}
    \Delta m = \bar{m}_{60} - \bar{m}_0
    \end{equation}
    Donde $\bar{m}_{60}$ es el promedio de los últimos 10 minutos del experimento.
    
    \item \textbf{Cálculo del Porcentaje de Reducción:}
    \begin{equation}
    \% \Delta m = \frac{\Delta m}{m_0} \times 100\%
    \end{equation}
    
    \item \textbf{Análisis de Temperatura:}
    \begin{equation}
    \Delta T = T_{60} - T_0
    \end{equation}
    \textbf{Criterio de éxito:} $\Delta T \leq 0^\circ$C (entropía negativa) o $\Delta T < +0.5^\circ$C (mínima disipación).
    
    \item \textbf{Análisis de Histéresis:}
    \begin{equation}
    \tau = \text{tiempo para que } \Delta m \text{ decaiga al 50\% tras apagar}
    \end{equation}
    Predicción del modelo: $\tau \approx 5-15$ minutos.
\end{enumerate}

% ==========================================
% SECCIÓN 4: CRITERIOS DE VALIDACIÓN
% ==========================================
\section{Criterios de Validación del Modelo GEM}

\begin{tcolorbox}[
    colback=gemgold!10,
    colframe=gemgold,
    title=\textbf{\textcolor{gemdark}{LAS 3 FIRMAS DEL ÉXITO (CRITERIOS DE VALIDACIÓN)}},
    fonttitle=\bfseries\large,
    boxrule=1.5pt,
    arc=4pt
]
Para considerar el experimento como \textbf{validación del Teorema 02}, deben cumplirse \textbf{simultáneamente} las siguientes condiciones:

\begin{enumerate}
    \item \textbf{Firma 1: Reducción de Masa Medible}
    \begin{itemize}
        \item $\Delta m > 3\sigma_0$ (señal mayor que 3 veces el ruido)
        \item $\% \Delta m \geq 0.01\%$ (reducción mínima del 0.01\%)
        \item Valor esperado: $0.01\% - 0.1\%$ ($\Delta m \approx 5-50$ mg para 50 g)
    \end{itemize}
    
    \item \textbf{Firma 2: Entropía Negativa o Neutra}
    \begin{itemize}
        \item $\Delta T \leq 0^\circ$C (enfriamiento = neguentropía) \textbf{IDEAL}
        \item $\Delta T < +0.5^\circ$C (mínima disipación) \textbf{ACEPTABLE}
        \item $\Delta T > +2.0^\circ$C \textbf{FALLO} (fuga vectorial, histéresis magnética)
    \end{itemize}
    
    \item \textbf{Firma 3: Histéresis Topológica}
    \begin{itemize}
        \item La reducción de masa persiste $> 5$ minutos tras apagar el campo
        \item Decaimiento exponencial con constante $\tau \approx 5-15$ min
        \item Esto demuestra que el efecto no es electromagnético transitorio, sino topológico
    \end{itemize}
\end{enumerate}
\end{tcolorbox}

\subsection{Interpretación de Resultados}

\begin{table}[h!]
\centering
\caption{Matriz de Interpretación de Resultados}
\begin{tabular}{@{}p{4cm}p{4cm}p{5cm}@{}}
\toprule
\textbf{Resultado} & \textbf{Interpretación} & \textbf{Acción} \\ \midrule
$\Delta m > 0$, $\Delta T \leq 0$, $\tau > 5$ min & \textbf{VALIDACIÓN EXITOSA} del Teorema 02 & Publicar datos, replicar en otro laboratorio \\ \addlinespace
$\Delta m > 0$, $\Delta T > 0$ (pequeño), $\tau > 5$ min & Validación parcial, posible fuga vectorial menor & Optimizar nulificación, repetir \\ \addlinespace
$\Delta m \approx 0$, $\Delta T \approx 0$ & Sin efecto detectable & Verificar calibración, aumentar amplitud \\ \addlinespace
$\Delta m < 0$ (aumento), $\Delta T > 2^\circ$C & \textbf{FALLO}: efecto Joule dominante & APAGAR, revisar bobinas, verificar fase \\ \bottomrule
\end{tabular}
\end{table}

% ==========================================
% SECCIÓN 5: PROTOCOLO DE SEGURIDAD
% ==========================================
\section{Protocolo de Seguridad y Contingencias}

\begin{tcolorbox}[
    colback=alertred!10,
    colframe=alertred,
    title=\textbf{\textcolor{alertred}{⚠️ REGLAS DE SEGURIDAD CRÍTICAS}},
    fonttitle=\bfseries,
    boxrule=1.5pt
]
\begin{enumerate}
    \item \textbf{Límite de Temperatura:} Si el termistor marca $\Delta T > +2.0^\circ$C, \textbf{APAGAR INMEDIATAMENTE} el generador. Esto indica fuga vectorial (histéresis magnética) que puede dañar las bobinas.
    
    \item \textbf{Límite de Corriente:} Monitorear corriente en bobinas. No exceder $I_{max} = 2$ A para evitar sobrecalentamiento.
    
    \item \textbf{Aislamiento Eléctrico:} Todas las conexiones deben estar aisladas. Usar guantes dieléctricos al manipular bobinas energizadas.
    
    \item \textbf{Campo Magnético:} Mantener objetos magnéticos (tarjetas, relojes, marcapasos) a > 1 m de distancia durante el experimento.
\end{enumerate}
\end{tcolorbox}

\subsection{Contingencias y Solución de Problemas}

\begin{table}[h!]
\centering
\caption{Tabla de Contingencias}
\begin{tabular}{@{}p{5cm}p{6cm}@{}}
\toprule
\textbf{Problema} & \textbf{Solución} \\ \midrule
La balanza muestra deriva térmica constante & Esperar más tiempo de estabilización (60 min). Verificar que la jaula de Faraday no tenga corrientes de aire. \\ \addlinespace
No se logra nulificación vectorial ($B > 0.1$ mT) & Verificar polaridad de bobinas. Ajustar fase con osciloscopio. Verificar que las bobinas sean perfectamente ortogonales. \\ \addlinespace
La muestra se calienta rápidamente ($\Delta T > 1^\circ$C/min) & APAGAR. Verificar que la frecuencia sea exactamente 16.2 Hz (no armónicos). Reducir amplitud a 2 Vpp y probar de nuevo. \\ \addlinespace
Ruido excesivo en la balanza ($\sigma > 0.5$ mg) & Verificar aislamiento de vibraciones (mesa óptica). Cerrar jaula de Faraday herméticamente. Desconectar equipos innecesarios. \\ \addlinespace
Efecto persistente tras apagar ($\tau > 30$ min) & Posible magnetización residual del vaso. Usar vaso de cuarzo en lugar de vidrio. Desmagnetizar bobinas con corriente alterna decreciente. \\ \bottomrule
\end{tabular}
\end{table}

% ==========================================
% SECCIÓN 6: REGISTRO DE DATOS
% ==========================================
\section{Formato de Registro de Datos}

\subsection{Hoja de Datos Experimental}

\begin{tcolorbox}[
    colback=gemblue!5,
    colframe=gemblue,
    title=\textbf{\textcolor{white}{HOJA DE REGISTRO GEM-01}},
    fonttitle=\bfseries
]
\textbf{Fecha:} \_\_\_\_\_\_\_\_\_\_\_\_\_\_\_\_\_\_\_\_ \quad \textbf{Operador:} \_\_\_\_\_\_\_\_\_\_\_\_\_\_\_\_\_\_\_\_ \\
\textbf{Temperatura Ambiente:} \_\_\_\_\_\_\_\_ $^\circ$C \quad \textbf{Humedad Relativa:} \_\_\_\_\_\_\_\_ \% \\[0.3cm]

\textbf{PARÁMETROS DEL EXPERIMENTO:} \\
Volumen de agua: \_\_\_\_\_\_\_\_ mL \quad Masa inicial $m_0$: \_\_\_\_\_\_\_\_ g \\
Frecuencia: \_\_\_\_\_\_\_\_ Hz \quad Amplitud: \_\_\_\_\_\_\_\_ Vpp \\
Campo magnético residual: $B_x = $\_\_\_\_\_, $B_y = $\_\_\_\_\_, $B_z = $\_\_\_\_\_ mT \\[0.3cm]

\textbf{RESULTADOS:} \\
Masa final ($t=60$ min): $\bar{m}_{60} = $\_\_\_\_\_\_\_\_ g \\
Reducción de masa: $\Delta m = $\_\_\_\_\_\_\_\_ g \\
Porcentaje: $\% \Delta m = $\_\_\_\_\_\_\_\_ \% \\
Temperatura final: $T_{60} = $\_\_\_\_\_\_\_\_ $^\circ$C \\
Cambio de temperatura: $\Delta T = $\_\_\_\_\_\_\_\_ $^\circ$C \\
Histéresis ($\tau$): \_\_\_\_\_\_\_\_ minutos \\[0.3cm]

\textbf{OBSERVACIONES:} \\
\rule{0.95\textwidth}{0.5pt} \\
\rule{0.95\textwidth}{0.5pt} \\
\rule{0.95\textwidth}{0.5pt} \\[0.3cm]

\textbf{FIRMA DEL OPERADOR:} \_\_\_\_\_\_\_\_\_\_\_\_\_\_\_\_\_\_\_\_ \quad \textbf{FECHA:} \_\_\_\_\_\_\_\_\_\_\_\_
\end{tcolorbox}

% ==========================================
% SECCIÓN 7: CONCLUSIÓN Y PRÓXIMOS PASOS
% ==========================================
\section{Conclusión y Próximos Pasos}

\subsection{Importancia de este Protocolo}

Este protocolo GEM-01 representa el \textbf{puente crítico} entre la teoría matemática del Modelo GEM y la realidad experimental. Si se validan las 3 firmas del éxito, habremos demostrado que:

\begin{enumerate}
    \item El vacío cuántico posee una presión base $P_0$ extraíble (Axioma 1 del Teorema 02).
    \item El gradiente topológico $\nabla w$ actúa como motor termodinámico (Axioma 2).
    \item La entropía negativa $S_w < 0$ es real y medible (Axioma 3).
    \item La constante $1836$ es efectivamente el factor de conversión entre energía del vacío y reducción de inercia.
\end{enumerate}

\subsection{Próximos Pasos tras la Validación}

\begin{enumerate}
    \item \textbf{Replicación Independiente:} Enviar este protocolo a 3 laboratorios independientes (universidades o centros de investigación) para replicación doble-ciego.
    
    \item \textbf{Publicación en Preprint:} Subir resultados a arXiv (physics.gen-ph) con todos los datos brutos y código de análisis.
    
    \item \textbf{Optimización del Rig:} Si $\Delta m$ es detectable pero pequeño, optimizar:
    \begin{itemize}
        \item Aumentar amplitud a 10 Vpp (si no hay calentamiento)
        \item Probar frecuencias cercanas: 15.8 Hz, 16.5 Hz (barrido fino)
        \item Añadir agua estructurada (vortexizado) para aumentar $C_{105}$
    \end{itemize}
    
    \item \textbf{GEM-02: Generación de Energía:} Si GEM-01 es exitoso, diseñar el siguiente experimento: conectar el sistema a un circuito de cosecha (diodos Schottky + condensador) para medir potencia eléctrica extraída.
\end{enumerate}

\begin{center}
\vspace{1cm}
\textbf{\Large "El territorio siempre tiene la última palabra."} \\
\vspace{0.3cm}
\textit{— Principio Fundamental del Modelo GEM}
\end{center}

\vspace{1cm}
\begin{center}
\textbf{Fin del Protocolo de Validación Experimental GEM-01} \\
Versión 1.0 — Junio 2026 \\
Licencia: CC BY-SA 4.0 (Ciencia Abierta)
\end{center}

\end{document}